\documentclass[conference]{IEEEtran}

\usepackage{cite}
\usepackage{amsmath,amssymb,amsfonts}
\usepackage{graphicx}
\usepackage{booktabs}
\usepackage{array}
\usepackage{url}
\usepackage{xcolor}

\title{Notes on the Gaussian Distribution}

\author{
\IEEEauthorblockN{Luis Mendez}
}

\begin{document}
\maketitle

\section*{Overview}
The Bayesian optimization notes lean on the Gaussian at three separate points without defining it: the surrogate is a \textit{Gaussian} process, the posterior stays closed-form because a Gaussian prior with a Gaussian likelihood is \textit{conjugate}, and Expected Improvement is written with $\Phi$ and $\phi$, the standard normal CDF and PDF. These notes fill in that gap. The goal is not a general probability review but the specific set of Gaussian facts that make the SMBO loop computable.

Three properties do essentially all the work:
\begin{itemize}
    \item \textbf{Closure under linear operations} --- affine maps and sums of independent Gaussians are Gaussian.
    \item \textbf{Closure under conditioning} --- marginals and conditionals of a joint Gaussian are Gaussian, with means and covariances given by block formulas.
    \item \textbf{Self-conjugacy} --- a Gaussian prior updated by a Gaussian likelihood gives a Gaussian posterior.
\end{itemize}
Every closed-form expression in the Gaussian process section is one of these three applied to a particular partition of variables. Nothing else about the distribution is needed to derive them.

\section*{The Univariate Gaussian}
A random variable $X$ is Gaussian (normal) with mean $\mu$ and variance $\sigma^2 > 0$, written $X \sim \mathcal{N}(\mu, \sigma^2)$, if its density is
\begin{equation}
    p(x) = \frac{1}{\sqrt{2\pi\sigma^2}}\exp\!\left(-\frac{(x-\mu)^2}{2\sigma^2}\right).
\end{equation}
The exponent is a negative squared distance from $\mu$, measured in units of $\sigma$; the prefactor is the constant that makes the density integrate to one. Two parameters fix the distribution completely, which is unusual --- there is no separate skewness or tail-weight knob to set.

The \textbf{standard} normal is the case $\mu = 0$, $\sigma^2 = 1$, whose density and CDF are written
\begin{equation}
    \phi(z) = \frac{1}{\sqrt{2\pi}}e^{-z^2/2},
    \qquad
    \Phi(z) = \int_{-\infty}^{z}\phi(t)\,dt.
\end{equation}
Any Gaussian reduces to it by standardization: if $X \sim \mathcal{N}(\mu,\sigma^2)$ then $Z = (X-\mu)/\sigma \sim \mathcal{N}(0,1)$. This is why implementations only ever need $\phi$ and $\Phi$, and why the $Z$ in the EI formula is a difference divided by a standard deviation.

$\Phi$ has no elementary closed form; it is evaluated numerically through the error function, $\Phi(z) = \tfrac{1}{2}\left[1 + \mathrm{erf}(z/\sqrt{2})\right]$. Moments:
\begin{equation}
    \mathbb{E}[X] = \mu, \qquad \mathrm{Var}[X] = \sigma^2, \qquad \mathbb{E}\!\left[(X-\mu)^3\right] = 0,
\end{equation}
so the distribution is symmetric about its mean, and mean, median, and mode coincide. Mass concentrates quickly: roughly $68\%$, $95\%$, and $99.7\%$ of it lies within one, two, and three standard deviations of $\mu$. That concentration is what makes $\sigma(\mathbf{x})$ usable as a confidence width in LCB --- $\kappa = 2$ is approximately a $95\%$ one-sided claim.

\subsection*{An aside on $\Phi$ versus the sigmoid}
$\Phi$ is a squashing function from $\mathbb{R}$ to $(0,1)$ with the same S-shape as the logistic sigmoid $\sigma(x) = 1/(1+e^{-x})$, and using it as the link in a binary classifier gives \textit{probit} regression rather than logistic regression. The two give very similar fits after rescaling; the logistic is preferred in practice because it is elementary, has the derivative identity $\sigma' = \sigma(1-\sigma)$, and pairs analytically with binary cross-entropy. The probit link is the version that arises when the latent variable driving the label is assumed Gaussian.

\section*{Closure Under Linear Operations}
If $X \sim \mathcal{N}(\mu, \sigma^2)$ and $Y = aX + b$ with $a \neq 0$, then
\begin{equation}
    Y \sim \mathcal{N}\!\left(a\mu + b,\; a^2\sigma^2\right).
\end{equation}
If $X_1 \sim \mathcal{N}(\mu_1,\sigma_1^2)$ and $X_2 \sim \mathcal{N}(\mu_2,\sigma_2^2)$ are independent, then
\begin{equation}
    X_1 + X_2 \sim \mathcal{N}\!\left(\mu_1 + \mu_2,\; \sigma_1^2 + \sigma_2^2\right).
\end{equation}
Note that variances add, not standard deviations. The consequence for modeling is direct: adding independent Gaussian observation noise to a Gaussian latent value keeps everything Gaussian and merely inflates the variance. That is exactly the $\sigma_n^2 I$ term added to the kernel matrix in the GP equations --- resampling noise in the cross-validated score, folded in without leaving the family.

\section*{The Multivariate Gaussian}
For $\mathbf{x} \in \mathbb{R}^d$ with mean vector $\boldsymbol{\mu}$ and symmetric positive-definite covariance $\Sigma$,
\begin{equation}
    p(\mathbf{x}) = \frac{1}{(2\pi)^{d/2}|\Sigma|^{1/2}}
    \exp\!\left(-\tfrac{1}{2}(\mathbf{x}-\boldsymbol{\mu})^{\!\top}\Sigma^{-1}(\mathbf{x}-\boldsymbol{\mu})\right).
\end{equation}
The quadratic form in the exponent is the squared \textbf{Mahalanobis distance}, the natural generalization of ``how many standard deviations away'' to correlated coordinates. Level sets of the density are therefore ellipsoids of constant Mahalanobis distance. Writing the eigendecomposition $\Sigma = U\Lambda U^{\!\top}$, those ellipsoids have axes along the eigenvectors (columns of $U$) with half-lengths proportional to $\sqrt{\lambda_i}$: the eigenvectors give the directions of correlated variation and the eigenvalues give the variance along each. Special cases are worth naming --- $\Sigma = \sigma^2 I$ gives spherical contours, a diagonal $\Sigma$ gives axis-aligned ellipsoids, and off-diagonal entries are what tilt them.

Affine closure carries over. For any matrix $A$ and vector $\mathbf{b}$,
\begin{equation}
    \mathbf{x} \sim \mathcal{N}(\boldsymbol{\mu},\Sigma)
    \;\Longrightarrow\;
    A\mathbf{x} + \mathbf{b} \sim \mathcal{N}\!\left(A\boldsymbol{\mu} + \mathbf{b},\; A\Sigma A^{\!\top}\right).
\end{equation}
Choosing $A = \Lambda^{-1/2}U^{\!\top}$ maps any Gaussian to a standard one, which is what \textit{whitening} means and how samples are generated: draw $\mathbf{z} \sim \mathcal{N}(\mathbf{0},I)$ and return $\boldsymbol{\mu} + L\mathbf{z}$, where $LL^{\!\top} = \Sigma$.

Two cautions that matter when reasoning about these objects. First, for jointly Gaussian variables, zero correlation \textit{does} imply independence --- a special property, not a general fact about random variables. Second, the converse of jointness fails: each coordinate being marginally Gaussian does not make the vector jointly Gaussian, so ``Gaussian'' applied to a collection is a genuinely stronger claim than it is applied to each member.

\section*{Marginals and Conditionals}
This is the section the Gaussian process rests on. Partition a jointly Gaussian vector and its parameters conformably:
\begin{equation}
    \begin{bmatrix}\mathbf{x}_1\\\mathbf{x}_2\end{bmatrix}
    \sim
    \mathcal{N}\!\left(
    \begin{bmatrix}\boldsymbol{\mu}_1\\\boldsymbol{\mu}_2\end{bmatrix},
    \begin{bmatrix}\Sigma_{11} & \Sigma_{12}\\ \Sigma_{21} & \Sigma_{22}\end{bmatrix}
    \right).
\end{equation}
The \textbf{marginal} is read straight off the blocks:
\begin{equation}
    \mathbf{x}_1 \sim \mathcal{N}(\boldsymbol{\mu}_1, \Sigma_{11}).
\end{equation}
Marginalizing out variables costs nothing --- delete their rows and columns. This is what licenses the GP's definition over an infinite index set while every computation touches only the finitely many points actually in play.

The \textbf{conditional} is the substantive formula:
\begin{align}
    \mathbf{x}_1 \mid \mathbf{x}_2 &\sim \mathcal{N}\!\left(\boldsymbol{\mu}_{1|2},\, \Sigma_{1|2}\right), \\
    \boldsymbol{\mu}_{1|2} &= \boldsymbol{\mu}_1 + \Sigma_{12}\Sigma_{22}^{-1}(\mathbf{x}_2 - \boldsymbol{\mu}_2), \\
    \Sigma_{1|2} &= \Sigma_{11} - \Sigma_{12}\Sigma_{22}^{-1}\Sigma_{21}.
\end{align}
Three things to read out of it. The conditional mean is the prior mean corrected by a linear function of how far the observation fell from \textit{its} prior mean, with $\Sigma_{12}\Sigma_{22}^{-1}$ playing the role of a regression coefficient. The conditional covariance is the prior covariance minus a positive-semidefinite term, so conditioning never increases uncertainty. And --- the key point --- $\Sigma_{1|2}$ does not depend on the observed value $\mathbf{x}_2$ at all, only on which variables were observed.

That last observation is the formal content of the claim in the Bayesian optimization notes that the GP's posterior variance depends on \textit{where} we sampled and not on the values obtained. It is a property of the Gaussian family, not of the kernel or the acquisition function, and it is what makes $\sigma(\mathbf{x})$ a clean exploration signal: uncertainty collapses wherever the search has looked, irrespective of whether the news was good or bad.

\section*{Conjugacy: The Posterior Update}
Take the simplest nontrivial case --- a Gaussian likelihood with known variance $\sigma^2$ and an unknown mean $\mu$ carrying a Gaussian prior:
\begin{equation}
    \mu \sim \mathcal{N}(\mu_0, \tau_0^2),
    \qquad
    x_i \mid \mu \sim \mathcal{N}(\mu, \sigma^2), \quad i = 1,\dots,n.
\end{equation}
Multiplying prior and likelihood gives a quadratic in $\mu$ inside an exponential, which is again a Gaussian density. Working in \textbf{precisions} (inverse variances) makes the result readable. With $\bar{x}$ the sample mean,
\begin{align}
    \frac{1}{\tau_n^2} &= \frac{1}{\tau_0^2} + \frac{n}{\sigma^2}, \\
    \mu_n &= \tau_n^2\left(\frac{\mu_0}{\tau_0^2} + \frac{n\bar{x}}{\sigma^2}\right),
\end{align}
so $\mu \mid x_{1:n} \sim \mathcal{N}(\mu_n, \tau_n^2)$. Precisions add, and the posterior mean is a precision-weighted average of the prior mean and the data mean. The two limits behave as they should: a vague prior ($\tau_0^2 \to \infty$) leaves $\mu_n \to \bar{x}$, pure data; and $n \to \infty$ drives $\tau_n^2 \to 0$, with the prior's influence vanishing at rate $1/n$.

This is the mechanism behind the remark that the GP escapes the intractable evidence integral. In general, Bayes' rule requires the normalizer $p(\mathcal{D}) = \int p(\mathcal{D}\mid\theta)p(\theta)\,d\theta$, which needs MCMC or a variational approximation. Under conjugacy the integral is a Gaussian integral, known analytically, so the update reduces to arithmetic on means and covariances. That is the difference between a surrogate that can be refitted inside every iteration and one that cannot.

\section*{From the Conditional Formula to the GP Posterior}
The GP equations quoted in the Bayesian optimization notes are now just the conditional formula applied once. A GP asserts that the function values at any finite set of inputs are jointly Gaussian, so with $X$ the observed inputs, $\mathbf{y}$ the observed (noisy) scores, $K$ the kernel matrix $K_{ij} = k(\mathbf{x}_i,\mathbf{x}_j)$, and $f(\mathbf{x}_*)$ the value at a new point, a zero-mean prior gives
\begin{equation}
    \begin{bmatrix}\mathbf{y}\\ f(\mathbf{x}_*)\end{bmatrix}
    \sim
    \mathcal{N}\!\left(\mathbf{0},
    \begin{bmatrix}
    K + \sigma_n^2 I & k(X,\mathbf{x}_*)\\
    k(\mathbf{x}_*,X) & k(\mathbf{x}_*,\mathbf{x}_*)
    \end{bmatrix}\right).
\end{equation}
The $\sigma_n^2 I$ block is independent observation noise added to each observed value --- the closure-under-sums property from earlier. Substituting the blocks into the conditional formulas, with $\mathbf{x}_1 = f(\mathbf{x}_*)$ and $\mathbf{x}_2 = \mathbf{y}$, gives
\begin{align}
    \mu(\mathbf{x}_*)      &= k(\mathbf{x}_*, X)\left[K + \sigma_n^2 I\right]^{-1}\mathbf{y}, \\
    \sigma^2(\mathbf{x}_*) &= k(\mathbf{x}_*, \mathbf{x}_*) - k(\mathbf{x}_*, X)\left[K + \sigma_n^2 I\right]^{-1}k(X, \mathbf{x}_*),
\end{align}
which are exactly the surrogate equations. The interpretation follows from the general formula: the mean is a similarity-weighted combination of observed scores because $\Sigma_{12}\Sigma_{22}^{-1}$ is a regression coefficient, and the variance subtracts a nonnegative term that is larger where the kernel says the new point resembles points already seen.

The $[K + \sigma_n^2 I]^{-1}$ inverse is also where the cubic scaling comes from. It is never formed explicitly; the standard route is a Cholesky factorization $K + \sigma_n^2 I = LL^{\!\top}$ followed by triangular solves, at $O(n^3)$ for the factorization and $O(n^2)$ per prediction. That cost is trivial at the tens-to-hundreds of trials a tuning budget allows and prohibitive at tens of thousands, which is the practical boundary noted among the GP's drawbacks.

\section*{Maximum Likelihood, and Why Gaussian Noise Means Least Squares}
For $n$ i.i.d.\ observations, maximizing the Gaussian likelihood gives
\begin{equation}
    \hat{\mu} = \frac{1}{n}\sum_{i=1}^{n}x_i,
    \qquad
    \hat{\sigma}^2 = \frac{1}{n}\sum_{i=1}^{n}(x_i - \hat{\mu})^2.
\end{equation}
The variance estimate is biased low by a factor $(n-1)/n$, since it measures spread about the estimated rather than the true mean; dividing by $n-1$ instead is the unbiased correction.

More consequential is what the Gaussian assumption implies about the loss function. Suppose $y_i = f(\mathbf{x}_i) + \varepsilon_i$ with $\varepsilon_i \sim \mathcal{N}(0,\sigma^2)$ independent. The negative log-likelihood is
\begin{equation}
    -\log p(\mathbf{y} \mid f) = \frac{n}{2}\log(2\pi\sigma^2) + \frac{1}{2\sigma^2}\sum_{i=1}^{n}\left(y_i - f(\mathbf{x}_i)\right)^2.
\end{equation}
Only the second term depends on $f$, and $\sigma^2$ scales it without moving the minimum, so \textbf{maximizing a Gaussian likelihood is minimizing squared error}. Least squares and mean squared error are not arbitrary choices of convenience --- they are the maximum-likelihood estimator under additive Gaussian noise, exactly as binary cross-entropy is the maximum-likelihood loss under a Bernoulli model with a sigmoid link. The choice of noise model \textit{is} the choice of loss. This also explains the pairing of the squared exponent with the quadratic penalty throughout: a Gaussian prior on parameters contributes $\|\mathbf{w}\|^2/2\tau^2$ to the same objective, which is ridge ($L_2$) regularization.

\section*{Why the Gaussian Keeps Appearing}
Three independent reasons, beyond analytical convenience:
\begin{itemize}
    \item \textbf{The central limit theorem}. Sums of many independent contributions with finite variance tend to a Gaussian regardless of the individual distributions. A cross-validated score is an average over folds and over many examples, so the resampling noise $\sigma_n^2$ in the GP is Gaussian for a real reason, not just a convenient one.
    \item \textbf{Maximum entropy}. Among all distributions on $\mathbb{R}$ with a given mean and variance, the Gaussian has the largest differential entropy, $\tfrac{1}{2}\log(2\pi e\sigma^2)$. Choosing it is therefore the least-committal choice consistent with knowing only a center and a spread --- the right default for a prior meant to assume as little as possible.
    \item \textbf{Closure}. As shown above, the family survives affine maps, sums, marginalization, conditioning, and Bayesian updating. No other continuous family is this stable, and that stability is what turns inference into linear algebra.
\end{itemize}

\section*{Practical Notes}
\begin{itemize}
    \item \textbf{Never invert; factorize}. Use Cholesky and triangular solves rather than an explicit $\Sigma^{-1}$: it is faster, better conditioned, and yields the log-determinant for free as $\log|\Sigma| = 2\sum_i \log L_{ii}$, which the marginal likelihood needs.
    \item \textbf{Add jitter}. A kernel matrix from nearly duplicate inputs is numerically singular. Adding a small $\varepsilon I$ (physically, the noise term $\sigma_n^2 I$) restores positive definiteness --- one reason to fit a noise level even when the objective is believed deterministic.
    \item \textbf{Work in log space}. Densities underflow in more than a few dimensions. Compute log-densities and use $\log$-$\sum$-$\exp$ rather than multiplying probabilities.
    \item \textbf{Standardize the targets}. GP implementations assume a zero-mean prior and length scales on a comparable scale across dimensions, so centering $\mathbf{y}$ and putting skewed hyperparameter ranges on a log scale both materially improve the fit.
    \item \textbf{The tails are a real assumption}. The Gaussian puts negligible mass beyond a few $\sigma$ and unbounded support on both sides, so a GP over an accuracy score can return a posterior mean above $1$ or a symmetric interval where the true distribution is skewed and bounded. This is usually harmless for ranking candidate points but is the reason heavier-tailed or warped alternatives exist, and it is worth remembering that a $2\sigma$ interval is a claim about the model, not a guarantee about the objective.
\end{itemize}

\end{document}
